% !TeX root = main.tex
\chapter{新型内存体系下数据管理机制的统一分析}

\section{三项工作的统一框架}
从研究层次、系统范围、关键硬件因素、核心问题、主要方法和优化目标六个维度对 ReCoW、MiTDM 和 Duet 进行对照。
\section{数据管理问题的层次演进}
\subsection{数据在哪里：动态位置的低开销表示}
\subsection{数据如何组织：远程访问下的布局优化}
\subsection{数据由谁管理：异构资源下的功能放置}
\section{硬件演进与研究范围扩展}
分析从持久内存、远程内存到计算—内存—存储解耦的系统边界变化，以及各章问题之间的递进关系。
\section{跨机制的通用设计原则}
\subsection{复用硬件已有机制}
\subsection{聚合昂贵远程访问}
\subsection{自适应访问粒度}
\subsection{按资源属性拆分功能}
\subsection{区分核心状态与辅助状态}
\section{适用范围与机制选择}
按照持久性需求、网络条件、工作负载、版本分布、容量压力和资源非对称程度讨论选择依据。
\section{局限性与潜在协同}
讨论页粒度重映射、自适应开销、资源假设，以及三项机制尚未整合为统一数据库运行时的局限。
\section{本章小结}
说明本章不是重复前三章，而是抽象跨层规律并建立完整论文论证闭环。
